\section{Retained buffered-TLS and fixed-position retail studies}\label{app:retained-controlled}
This section preserves, unchanged in substance, the first-revision main-text account of the two controlled studies that Section~8 of the main paper now summarizes. Section~\ref{app:transfer} reports the second-round tests E5--E7.

The two studies test distinct observers. The transport observer sees TLS application-record lengths on a selected link. The tool-dispatch observer sees ordered tool names, write counts and response byte lengths, including failed calls; it is an instrumented dispatch view, not an encrypted-network observer. Neither experiment combines these views or runs an autonomous agent. Saved outputs and reference programs let us vary the specific boundary while holding the prescribed workload fixed.

\subsection{Transport placement and workload-matched framing}
For 24 public topics, the saved corpus contains one prompted brief and one prompted expanded \texttt{gpt-6-luna} response (48 outputs; 12 training and 12 disjoint test topics). These are designed prompt conditions, not a sample of natural task lengths. The original placement experiment passes text through two TLS~1.2 connections: the frontend decodes and removes upstream framing before sending new terminal chunks. Exact plaintext hashes are checked. Equalizing the upstream length projection leaves the unpadded terminal projection separable (100\% verbosity classification); padding both projections reduces the evaluated classifier to 50\%. This reproduces the receiving-boundary mechanism identified by NetEcho rather than discovering a new failure mode.

The revised framing replay uses \texttt{ECDHE-RSA-AES128-GCM-SHA256} sockets and parses application-data records. T1 contrasts brief versus full expanded responses. T2 contrasts the brief with the first two period-delimited sentences of the expanded response. T3 contrasts the first two versus first three sentences of the same expanded response. Thus T2/T3 are derived prefixes, with two examples per topic, not independently generated explanations. T4 relabels both brief and expanded responses by topic domain; it has 24 test outputs from 12 topics, not 12 expanded outputs. Its unpadded accuracy is 45.8\% (95\% topic-bootstrap interval 25.0--66.7\%); its test majority-class baseline is 58.3\%. T4 is a weak-signal control, not a defense success.

\begin{table}[!t]\centering\footnotesize
\caption{Corrected TLS application-record results. Costs use each workload's own unprotected texts (all 48 outputs); accuracy uses 24 test outputs in 12 topics. Intervals resample topics.}
\label{tab:controlled-placement}
\setlength{\tabcolsep}{2.2pt}
\begin{tabularx}{\columnwidth}{@{}Yrrr@{}}\toprule
\multicolumn{4}{@{}l}{\textbf{A. Two-link placement (T1; retained experiment)}}\\
Protection & Upstream acc. & Terminal acc. & Terminal bytes\\\midrule
None & 100.0 & 100.0 & 720.96\\
Upstream only & 50.0 & 100.0 & 720.96\\
Both, loose cap & 50.0 & 50.0 & 16,848\\\midrule
\multicolumn{4}{@{}l}{\textbf{B. Matched workload costs and defended-fit accuracy}}\\
Workload / framing & Mean bytes & Ratio & Acc. (\%)\\\midrule
T1: unprotected & 720.96 & 1.00 & 100.0\\
T1: 256B bucket & 856.96 & 1.19 & 100.0\\
T1: 512B bucket & 1,085.92 & 1.51 & 100.0\\
T1: 1024B bucket & 1,401.25 & 1.94 & 70.8\\
T1: training-fitted cap & 1,570.00 & 2.18 & 50.0\\
T1: loose cap & 16,848.00 & 23.37 & 50.0\\
T2: unprotected & 272.48 & 1.00 & 100.0\\
T2: 256B bucket & 438.29 & 1.61 & 87.5\\
T2: 512B bucket & 628.00 & 2.30 & 50.0\\
T2: 1024B bucket & 1,140.00 & 4.18 & 50.0\\
T3: unprotected & 456.33 & 1.00 & 91.7\\
T3: 256B bucket & 608.38 & 1.33 & 62.5\\
T3: 512B bucket & 654.17 & 1.43 & 50.0\\
T3: 1024B bucket & 1,140.00 & 2.50 & 50.0\\
T4: unprotected & 720.96 & 1.00 & 45.8\\
\midrule
\multicolumn{4}{@{}l}{\textbf{C. Adaptation on identical defended T1 tests}}\\
Defense & $D_0$ (\%) & $D_2$ (\%) & Gain [95\% CI], pp\\\midrule
+512B offset & 50.0 & 100.0 & 50.0 [50.0, 50.0]\\
1024B bucket & 50.0 & 70.8 & 20.8 [8.3, 33.3]\\
$U(0,512)$B & 86.4 & 99.9 & 13.5 [10.7, 16.2]\\
$U(0,1024)$B & 68.2 & 89.2 & 21.1 [17.2, 24.8]\\
Loose cap & 50.0 & 50.0 & 0.0 [0.0, 0.0]\\
Training-fitted cap & 50.0 & 50.0 & 0.0 [0.0, 0.0]\\
\bottomrule\end{tabularx}
\par\smallskip\raggedright\footnotesize
Panel C uses 12 training pairs. Four-pair accuracy is the same except under $U(0,1024)$B, where it is 90.2\%. Random padding averages 50 draws within each topic; these are not 50 independent corpora. T2/T3 are sentence-prefix constructions. T4 relabels both original response lengths and has a 58.3\% test majority baseline. The cap equalizes only the specified record-length projection when no message overflows.
\end{table}


Every cost ratio divides defended application-record bytes by unpadded application-record bytes for the \emph{same texts}, over all 48 examples in that workload. The count includes TLS record headers and encryption overhead; handshakes, TCP/IP, upstream bytes and latency are excluded. The 4-byte length prefix is included. Buckets of 256/512 bytes use 128-byte writes; 1024-byte buckets use 256-byte writes. Consequently the 1024-byte single-quantum wire size is 1,140 bytes, not 1,082. The workload-matched 512-byte ratios for T2/T3 are 2.30$\times$/1.43$\times$; dividing both by T1's 720.96-byte denominator would yield the invalid earlier ratios below one.

A tighter comparison selects a constant application-frame cap using \emph{training lengths only}: round the largest framed training response (1,194 bytes) up to a multiple of 128, yielding 1,280 bytes. Ten 128-byte records cost 1,570 wire bytes (2.18$\times$ T1). Longer messages would emit additional 128-byte records without truncation; no test message overflows. This produces identical length sequences on the saved T1 test set and 50\% accuracy for the evaluated classifiers. It is a feasible cap for these data, not a globally optimal padding protocol or a guarantee for future lengths. The retained 16,848-byte cap (23.37$\times$) is only a deliberately loose control.

\subsection{Matched attacker adaptation}
We use $D_0$ for an attacker fitted on unprotected training data and $D_2$ for one fitted on defended training data; both are evaluated on the identical defended test examples. The threshold uses total TLS record payload bytes. $D_2$ uses the first four or all twelve training topic pairs; these budgets do not represent multiple independently sampled training sets. At a 1024-byte bucket, $D_0$ constantly predicts the expanded class and scores 50\%, while $D_2$ scores 70.8\% (58.3--83.3\%). An offset of 512 bytes permits 100\% adapted accuracy. This demonstrates why a shifted $D_0$ alone is insufficient, without claiming novelty for adaptive evaluation.

Random-tail padding is evaluated over 50 draws per output. The first draw is socket measured and the remaining 49 use the record-length accounting checked against the socket replay. Each topic's paired correctness is averaged over its outputs and draws before 10,000 topic-bootstrap resamples (seed 20261021). Fractional accuracies therefore do not imply additional independent topics; the intervals condition on these saved outputs and fitted attackers. Positive adaptation gains span 13.5--50.0 percentage points; both constant caps have zero gain.

No measured accuracy of 50\% proves absence of all leakage. The constant-cap construction makes only the specified record-length sequence constant when there is no overflow. Timing, connection behavior, request metadata and content are outside that projection. No sequence/timing attack, TLS~1.3 native padding, QUIC or live token/SSE stream is evaluated.

\subsection{Retail execution: state scope and a visible-call counterexample}\label{app:retained-retail}
We replay the 115 public reference programs in the pinned $\tau$-bench retail simulator. There is no LLM planner or user simulator in this experiment. Ninety-seven tasks have an initially pending order owned by the task user; 18 have none and receive no cover. Three seeds select orders, and cover is inserted before or after the prescribed program. Utility uses 97 task clusters, each with three repeated units per position. Policies include no cover, discard, shadow execution, read only, address self-write, item self-write, instruction-aware selection, other-account cancellation and blind cancellation. Destructive conditions are local negative controls, not proposed deployment policies.

\paragraph{Separate the outcome scopes}
Table~\ref{tab:tau-cover-bridge} reports reference database equality, full active-account equality and equality on objects \emph{changed} by the reference program. The last scope does not include every read/accessed object. The database criterion corresponds to the state part of the native evaluator; the extension script does not evaluate generated response correctness and we do not relabel its state equality as an end-to-end native task reward. All-account equality now checks every order owned by the active user and that user's record, rather than only initially pending orders.
\begin{table}[t]\centering\footnotesize
\caption{Retail state and observer results. State columns use 97 eligible tasks (291 units per position); observer columns use 49 held-out tasks (147 units). Percentages are task means across three seeds. B/A denotes before/after.}\label{tab:tau-cover-bridge}
\setlength{\tabcolsep}{2.2pt}
\begin{tabularx}{\columnwidth}{@{}Yrrrr@{}}\toprule
\multicolumn{5}{@{}l}{\textbf{A. Distinct state scopes (reference execution as comparator)}}\\
Policy / position & DB eq. & Account eq. & Object eq. & Mutated\\\midrule
No cover B/A & 100.00 & 100.00 & 100.00 & 0.00\\
Discard B/A & 100.00 & 100.00 & 100.00 & 0.00\\
Shadow B/A & 100.00 & 100.00 & 100.00 & 0.00\\
Read only B/A & 100.00 & 100.00 & 100.00 & 0.00\\
Address self-write B/A & 100.00 & 100.00 & 100.00 & 0.00\\
Item self-write B & 40.55 & 40.55 & 70.10 & 59.45\\
Item self-write A & 67.70 & 67.70 & 97.25 & 32.30\\
Instruction aware B & 35.05 & 35.05 & 78.35 & 70.10\\
Instruction aware A & 46.39 & 46.39 & 89.69 & 53.61\\
Other-account cancel B & 0.00 & 100.00 & 100.00 & 100.00\\
Other-account cancel A & 0.00 & 100.00 & 100.00 & 100.00\\
Blind cancel B & 11.34 & 11.34 & 62.54 & 100.00\\
Blind cancel A & 41.24 & 41.24 & 92.44 & 58.76\\
\midrule
\multicolumn{5}{@{}l}{\textbf{B. Legacy learned observer versus allowed-view rules}}\\
Policy / position & $D_0$ acc. & $D_2$ acc. & $D_2$ BA & Rule acc.\\\midrule
Majority (all conditions) & 77.55 & 77.55 & 50.00 & ---\\
Instruction aware B & 51.02 & 73.47 & 47.37 & 100.00\\
Instruction aware A & 51.02 & 71.43 & 52.51 & 100.00\\
Other-account cancel B & 22.45 & 71.43 & 46.05 & 100.00\\
Other-account cancel A & 22.45 & 71.43 & 46.05 & 100.00\\
Blind cancel B & 22.45 & 71.43 & 46.05 & 100.00\\
Blind cancel A & 22.45 & 68.03 & 47.09 & 100.00\\
\bottomrule\end{tabularx}
\par\smallskip\raggedright
DB eq.: equality of orders/users against reference (products are read-only), the native reward's database criterion, not an output-response/authorization check. Account eq.: all orders owned by the active user plus that user's record. Object eq.: objects changed by the reference program, not all accessed objects. Mutated: cover itself changes state. Other-account calls yield 0\% DB equality but 100\% active-account equality. BA: balanced accuracy. Rule: remove the known first/last cover call, then inspect cancellation verbs; blind/other-account cancellation also admits the $N_{\rm cancel}\ge2$ rule. Both rules attain 100\% acc./BA in their stated conditions. Other non-cancellation policies retain 100\% verb inference. Full intervals, AUROC and confusion matrices appear in Supplement S9 and the JSON artifact.
\end{table}


Address self-writes preserve the reference state in this simulator. Item self-writes mutate state in 173/291 pre-insertion units (59.45\%): the tool changes the status to \texttt{pending (item modified)} and appends a refund entry even when the item identifier and price are unchanged. The remaining 118/291 calls fail because the requested variant is unavailable, explaining why this mutation rate is below blind cancellation's 291/291 successful calls. This is a state-preservation counterexample from one API, not a prevalence estimate or a discovery about idempotence. Formally, a no-change requirement $f(s,a)=s$ differs from repeated-call idempotence $f(f(s,a),a)=f(s,a)$.

Instruction-aware cover uses only instruction text and the initial order table for selection, with address self-write as fallback. It preserves reference-changed objects in 78.35\%/89.69\% of before/after units, but preserves the full account and database in only 35.05\%/46.39\%. Cancelling an unmentioned order can therefore damage a bystander resource. The heuristic was developed using the available task collection; it has no held-out policy-development evaluation. Its 68/97 availability (70.10\%) and the item-self-write object equality of 204/291 (70.10\%) are separate counts that happen to yield the same percentage.

The other-account condition selects a pending order belonging to another simulated user. It is not a provisioned dedicated decoy account. It preserves the active user's complete account and reference-changed objects in these tasks, while changing the shared database in every eligible unit (0\% database equality). Thus account isolation of an effect cannot be described as whole-database preservation or an authorized production defense.

\paragraph{Observer permissions and simple attacks}
The secret label is whether the reference program \emph{attempts} a cancellation, including a call that may fail after cover execution. The observer sees the ordered dispatched tool names, write counts and response lengths; failed calls remain visible. It sees no arguments, user/order identifiers or database state. We do not infer that a passive TLS observer has these permissions. A backend with plaintext state would instead need a different confidentiality target.

The fixed task split contains 48 training tasks (7 positive, 41 negative) and 49 test tasks (11 positive, 38 negative); seeds stay within tasks. Always predicting no cancellation gives 77.55\% test accuracy and 50\% balanced accuracy. The legacy $D_2$ linear classifier uses cancellation presence, write count and write-response bytes. Its 68.03--73.47\% accuracy under cancellation cover is below the majority baseline, and its balanced accuracy ranges from 46.05\% to 52.51\%. An increase over $D_0$ is therefore insufficient evidence for its claimed privacy interpretation.

For blind and other-account cancellation, exactly one extra cancellation is dispatched in every eligible session. The permitted rule $\widehat y=\mathbf{1}[N_{\rm cancel}\ge2]$ recovers the reference label at 100\% accuracy and balanced accuracy on all 49 test tasks and three seeds, before and after insertion. For any policy with one live cover call at a known first/last position, removing that call and testing for a remaining cancellation also gives 100\%. It needs no identifiers, no response content and no training labels. These are measured finite-sample results and a structural property of the fixed insertion design, not proof of general attack performance. Identifier filtering is unnecessary in this view and is not claimed as an evaluated attack.

\paragraph{Terminal equality does not establish path properties}
Blind pre-insertion cancellation leaves the final database equal to reference in 33/291 units across 12 tasks, even though it preempts the later reference cancellation. The exact per-task counts are in Supplement S9. After insertion, 120/291 cover calls are rejected because reference execution already changed the order's status, and those rejections explain the 41.24\% equality rate. Rejection is not successful benign cover. In the 82 tasks with no reference execution error (246 units), before/after equality is 24/246 (9.76\%) and 94/246 (38.21\%). These results illustrate the limits of endpoint scoring; consent and authorization were not measured.

